\documentclass[10pt,a4paper]{amsart}
\usepackage[margin=19mm]{geometry}
\usepackage{amsmath,amssymb,mathtools}
\usepackage[scr=rsfs,cal=euler]{mathalfa}
\usepackage{xfrac}
\usepackage[unicode,hidelinks,bookmarks=false]{hyperref}
\newcommand{\ie}{i.e.\ }
\newcommand{\cf}{cf.\ }
\newcommand{\resp}{respectively\ }
\newcommand{\Subsection}{\subsection}
\providecommand{\nobreakdash}{\nobreak\hbox{-}}
\setlength{\emergencystretch}{2em}
\multlinegap0pt
\usepackage{bm}
\usepackage{bbm}
\usepackage{dsfont}
\usepackage{xspace}
\usepackage{cancel}
\usepackage{mathtools}
\usepackage{amsthm}
\makeatletter
\@ifundefined{theorem}{\newtheorem{theorem}{Theorem}}{}
\@ifundefined{lemma}{\newtheorem{lemma}[theorem]{Lemma}}{}
\makeatother
\newcommand{\PR}[1]{{\mathbb{P}}\left\{ #1\right\}}
\newcommand{\cQ}{\ensuremath{\mathcal{Q}}}
\title[On the Error Probability of RPA Decoding of Reed-Muller Code]{On the Error Probability of RPA Decoding of Reed-Muller Codes over BMS Channels}
\author{Dorsa Fathollahi, V. Arvind Rameshwar, V. Lalitha}
\date{Source: arXiv:2601.09581v1 (CC BY 4.0)}
\begin{document}
\maketitle

\noindent\emph{Source and licence.} On the Error Probability of RPA Decoding of Reed-Muller Codes over BMS Channels. Version arXiv:2601.09581v1 (CC BY 4.0), distributed under the Creative Commons Attribution 4.0 International licence, as recorded in the arXiv licence metadata for this exact version and shown on the abstract page https://arxiv.org/abs/2601.09581v1. Full rights notice: licenses/arxiv-2601.09581-CC-BY-4.0.txt.\par\smallskip

\noindent\textit{Editorial scope.} Counted unit: the Lemma whose source label is \texttt{lem:Qi-unroll} in the article (source0.tex lines 659--682, source label \texttt{lem:Qi-unroll}), delivered together with the article's own complete original proof of it (source0.tex lines 700--767, one \texttt{proof} environment of 68 lines). No mathematics is written, repaired or re-derived: statement and proof are the article's own text line for line. The proof is detached from the statement in the source: the article states the result at lines 659--682 and prints its proof at lines 700--767, and the proof environment's own header names the statement, so the association is the article's own. Required context retained uncounted: lem:bhat-unroll (lines 485--488); thm:main (lines 497--507); lem:Q1 (lines 532--543). Every definition, display and label of those lines is retained unchanged. Omitted from this package and not redistributed: 1--484, 489--496, 508--531, 544--658, 683--699, 768--777. Nothing inside a retained excerpt is omitted or altered: every retained body is delivered exactly as the article's lines are written, so no omission receipt of any kind accompanies this package. Citations: the retained excerpts carry no citation command at all, so no bibliography entry of the article is transcribed and no reference is redistributed. Reference adaptation: none was needed and none was made; every reference inside the delivered unit resolves inside it, so no unresolved reference or citation is produced. Printed slips kept literal: no printed word, symbol, hypothesis or number of the counted statement or of its proof was altered. Layout and class: the article's own class is \texttt{IEEEtran} with packages including cite, amsmath, amssymb, amsfonts, graphicx, textcomp, xcolor, geometry, caption, hyperref, amsthm, enumitem, fullpage, nicefrac; the wrapper is the lane's pinned \texttt{amsart} editorial wrapper at 10pt on A4, which restates the theorem-like environments the retained text needs and adds the article's own macro definitions that the retained text needs (\texttt{\textbackslash PR}, \texttt{\textbackslash cQ}), with extra wrapper packages (bm, bbm, dsfont, xspace, cancel, mathtools). The delivered numbering is the unit's own. Assets: no retained span contains a figure environment, a graphics-inclusion command, a PDF-page inclusion, a table or a TikZ picture; no image file is copied into this package, the package declares no asset dependency and no image file is redistributed with it. Licence: version arXiv:2601.09581v1 of this article is distributed under the Creative Commons Attribution 4.0 International licence, as recorded in the arXiv licence metadata for this exact version and shown on the abstract page https://arxiv.org/abs/2601.09581v1. A full rights notice is delivered at licenses/arxiv-2601.09581-CC-BY-4.0.txt. Characters: every retained excerpt is pure ASCII, so the delivered engine, XeLaTeX with its default Latin Modern text font, reports no missing character.
\begin{lemma}\label{lem:bhat-unroll}
We have that for any $i\in \{1,\ldots,r\}$, $ Z_i  \leq 1- (1- Z_r)^{2^{r-i}}$.
    % Let $Z_r = Z(W)$ where $W$ is the original channel  that the codewords of  $RM(m,r)$ are sent over. 
\end{lemma}

\begin{theorem}\label{thm:main}
Let
$
\lambda := -\ln(1-Z)\in(0,\infty).
$
For
$
r < \log_2 m -\log_2 \lambda,
$
we have $P_\text{err}(\text{RM}(m,r))\xrightarrow{m\to \infty} 0$.
\end{theorem}

\begin{lemma}\label{lem:Q1}
     We have that
$
        \PR{ \cQ_1}  \leq (2 ^ { m -r + 2  } - 1 ) Z_1^{2^{m-r}}.
$
   %  Hence, it follows that
   % \begin{equation}
   %       \PR{ \overline { \cQ_1} } \geq 1 -  (2 ^ { m -r + 2  } - 1 ) Z_1^{2^{m-r}}. 
   % \end{equation}

   % here the $2^{m-r}$ in the power above $Z_1$ stems from $ d_{min} (RM(m-r+1, 1 ) ) =  2^{m-r} $.  
\end{lemma}

\begin{lemma}\label{lem:Qi-unroll}
Define for each $t\in\{2,\dots,r\}$,
\[
A_t := 2^{m-r+t}\, Z_t^{\,2^{m-r+t}-1},
\qquad
B_t := 2^{m-r+t}-1.
\]
Then for every $i\in\{2,\dots,r\}$,
\begin{equation*}\label{eq:Qi-unrolled}
\PR{\cQ_i}
\le
\sum_{t=2}^{i}\Bigg( A_t \prod_{s=t+1}^{i} B_s \Bigg)
+
\PR{\cQ_1}\prod_{s=2}^{i} B_s .
\end{equation*}
In particular, 
\begin{equation*}\label{eq:Qr-unrolled}
\PR{\cQ_r}
\le
\sum_{t=2}^{r}\Bigg( A_t \prod_{s=t+1}^{r} B_s \Bigg)
+
\PR{\cQ_1}\prod_{s=2}^{r} B_s .
\end{equation*}
\end{lemma}

\begin{proof}[Proof of Thm.~\ref{thm:main}]
    For each $ t \in  \{ 2 , \ldots, r \}$ from Lemma~\ref{lem:bhat-unroll} we have that $A_t := 2^{m-r+t}  Z_t^{2^{m-r+t}-1}$ and $B_t := 2^{m-r+t}-1 \leq 2^{m-r+t}.$
    % \begin{equation}
    %     A_t := 2^{m-r+t}  Z_t^{2^{m-r+t}-1},\ B_t := 2^{m-r+t}-1 \leq 2^{m-r+t}.  
    % \end{equation}
    % We now show that $A_t$ is very small due to its containing powers of $Z_t$ and that multiples of $B_t$ are negligible, besides. 
    Observe that for any $u\in \{2,\ldots,r\}$, we have
    \begin{align*}
       \prod_{s=u}^r B_s &=  \prod_{s=u}^r 2^{m-r+s} \\&=2^{(r-u+1)(m-r)+\sum_{s=u}^r s}
\leq 2^{O(mr)}.
    \end{align*}
Thus, from Lemma \ref{lem:Qi-unroll},
\begin{equation}
\label{eq:temp}
    \PR{\cQ_r}
\leq 2^{O(mr)} \PR{ \cQ_1} + \sum_{t=2 }^r 2^{ O(mr)} Z_t^{2^{m-r+t}-1}.
\end{equation}

We will individually show that each term above approaches zero in the large $m$ limit, for $r \le \log_2 m -\log_2 \lambda-\delta$, for some small constant $\delta$. For the first term we will show that $\PR{\cQ_1}$ is very small compared to its multiplicative pre-factor that is $2^{ O(mr)}$. To this end, recall that $ \lambda := -\ln( 1- Z) = -\ln( 1- Z_r)$. From Lemma~\ref{lem:bhat-unroll}, we know that
$
    ( 1 - Z_t ) \geq ( 1 - Z_r)^{2^{r-t}} = e^{-\lambda \cdot 2^{r-t}}.
$
Now, from Lemma~\ref{lem:Q1}, we obtain using the previous inequality that
\begin{equation*}\label{eq:asym-Q1}
\begin{aligned}
        \PR{ \cQ_1} % &\leq 2 ^ { m -r + 2  }\cdot Z_1^{2^{m-r}}\\
       & \leq  2^{m-r+2} e^{-(2^{m-r}\cdot  e^{  -\lambda\cdot 2^{r-1}})}.
       \end{aligned}
\end{equation*} 
Here, we use the fact that for any positive integer $k$, $ Z_t^{k} \leq e^{ -k ( 1- Z_t)  } $. Notice that in the regime $r \le \log_2 m -\log_2 \lambda -\delta$, we have $ \lambda\cdot 2^{r-1} = O(m)$. Therefore,
$
      2^{m-r}\cdot\exp(-\lambda 2^{r-1})=\exp(\Omega(m)).
$
Plugging this into the first term in \eqref{eq:temp}, we get that
\begin{equation*}
    \begin{aligned}
        2^{O(mr)} \PR{ \cQ_1}  &\leq 2^{O(mr)}  2^{m-r+2} e^{-(2^{m-r}  e^{  -\lambda 2^{r-1}})}\\
        & \leq 2^{O(mr)} e^{-e^{\Omega (m) }}\xrightarrow{m\to \infty} 0.
    \end{aligned}
\end{equation*}
% The second part being double exponential guarantees that, $2^{O(mr)} \PR{ \cQ_1} \rightarrow 0$ as $ m \rightarrow \infty$. 

Now we focus on the second summation in \eqref{eq:temp}; we will show that each summand approaches zero. So fix any $t$, define $b_t = 2^{ m-r + t } - 1 $. Then, since
$
    Z_t^{b_t} \leq e^{- b_t(1-Z_t) },
$
we have that
\begin{equation*}
\begin{aligned}  
        (1-Z_t ) b_t  &\geq \exp (-\lambda 2^{r-t} ) (  2^{m-r+t}  - 1 ) \\
        &\geq   2^{m-r+t-1 }\cdot \exp (-\lambda\cdot 2^{r-t} )\\
        &= \exp\left({ ( m - r + t -1 )\ln 2 - \lambda\cdot 2^{r-t}}\right).
 \end{aligned}
\end{equation*}
Again using the fact that $ \lambda\cdot 2^{r-t}\leq \lambda\cdot 2^{r-1} = o(m) $, we get that
\begin{equation*}
    ( m - r + t -1 )\ln 2 - \lambda\cdot 2^{r-t} = \Theta(m) - o(m) = \Omega(m). 
\end{equation*}
Plugging this into any term in the summation in \eqref{eq:temp}, we get that 
\begin{equation*}
    \begin{aligned}
        2^{ O(mr) } Z_t^{2^ { m-r+t} -1 }  &\leq 2^{O(mr) } e^{ -\exp( m - r + t -1 )\ln 2 - \lambda 2^{r-t})} \\ 
        % &= \exp( O (mr) \cdot \exp ( -\exp ( \Omega(m)) )\\ 
        &\xrightarrow{m\to \infty} 0.
    \end{aligned}
\end{equation*}
This concludes the proof of Theorem \ref{thm:main}.
\end{proof}

\end{document}
